% SUDS Showcase poster v8, August 14, 2026
% 4 ft x 3 ft landscape (121.92 cm x 91.44 cm), beamerposter.
% Build:  cd poster && pdflatex suds_poster_v8.tex
% Figures: every figure used by v8 is pinned in artifacts/figures_v8/.

\newcommand{\AuthorOne}{Yasir Alsugair}
\newcommand{\AuthorTwo}{Chuxuan Ai}

\documentclass[final]{beamer}
\usepackage[size=custom,width=121.92,height=91.44,scale=1.32]{beamerposter}
\usepackage[T1]{fontenc}
\usepackage{lmodern}
% big operators come from the extension font at its 10pt design size, so without
% this \int and \sum stay small while everything else scales up to poster size
\usepackage{exscale}
\renewcommand{\familydefault}{\sfdefault}
% beamer's default sans math has no sans Greek, so it silently mixes a serif
% theta into sans variables. One serif math family throughout instead.
\usefonttheme[onlymath]{serif}
\usepackage{graphicx}
\usepackage{amsmath}
\usepackage{qrcode}
\graphicspath{{artifacts/figures_v8/}{artifacts/figures_v7/}{artifacts/figures_v6/}{artifacts/figures/}{figures/}}

\setbeamertemplate{navigation symbols}{}
\setbeamersize{text margin left=2.5cm, text margin right=2.5cm}

\definecolor{uoftblue}{HTML}{002A5C}
\definecolor{inktext}{HTML}{20222B}
\definecolor{panelbg}{HTML}{F4F5F8}
\definecolor{softgray}{HTML}{8A8D99}
% Figure-local keys: the plots are allowed their own palettes, but nearby prose
% must repeat the colours visible in the corresponding figure.
\definecolor{gaiamap}{HTML}{9AA0AC}
\definecolor{gaiaensemble}{HTML}{3976B9}
\definecolor{gaiasghmc}{HTML}{8064A2}
\definecolor{gaiart}{HTML}{E6A11A}

\newcommand{\dd}{\mathrm{d}}

% numbered marker pointing into the References block at the foot of column 3
\newcommand{\refmark}[1]{\textsuperscript{\textcolor{uoftblue}{[#1]}}}
% one reference line, hung so the number stays in the left margin
\newcommand{\refitem}[2]{\noindent\hangindent=1.5em\hangafter=1 [#1]~#2\par}

% boxed key equations: \boxed inherits \fboxsep, whose 3pt default crowds a tall
% fraction against the rule at poster scale
\newcommand{\eqbox}[1]{{\setlength{\fboxsep}{9pt}\boxed{#1}}}

% Header band shares its outer edges with the block bars below, which run
% \bandmargin .. \paperwidth-\bandmargin on this geometry. Both the band and
% \linewidth are centred on the paper, so centring a \bandwidth box inside
% \linewidth lands on those edges without depending on beamer's column widths.
\newlength{\bandmargin}\setlength{\bandmargin}{2.13cm}
\newlength{\bandwidth}\setlength{\bandwidth}{\paperwidth}
\addtolength{\bandwidth}{-2\bandmargin}
\newcommand{\bandrule}{%
  \noindent\hbox to \linewidth{\hss
    \textcolor{uoftblue}{\rule{\bandwidth}{1.8mm}}%
  \hss}}

% Three logos on one line. The band gets an explicit height so what follows it
% does not depend on how deep the centred title box happens to be, and the
% flanking logo groups are set in zero-width boxes so the title stays centred on
% the paper whatever those groups measure.
\newlength{\bandheight}\setlength{\bandheight}{10.6cm}
\newcommand{\headerband}[1]{%
  \noindent\hbox to \linewidth{\hss
    \begin{minipage}[c][\bandheight][c]{\bandwidth}\leavevmode#1\end{minipage}%
  \hss}}

\setbeamercolor{normal text}{fg=inktext}
\setbeamercolor{block title}{fg=white,bg=uoftblue}
\setbeamercolor{block body}{fg=inktext,bg=panelbg}
\setbeamerfont{block title}{size=\Large,series=\bfseries}
\setbeamercolor{itemize item}{fg=uoftblue}

\setbeamertemplate{block begin}{%
  \begin{beamercolorbox}[ht=3.6ex,dp=1.2ex,leftskip=1.2ex]{block title}%
    \usebeamerfont*{block title}\insertblocktitle
  \end{beamercolorbox}%
  \nointerlineskip
  \begin{beamercolorbox}[vmode]{block body}%
    \vspace{0.9ex}%
    % Use a genuinely narrower content area so text, figures, and side-by-side
    % minipages all share the same small inset and see the correct \linewidth.
    \noindent\hspace*{1.2ex}%
    \begin{minipage}{\dimexpr\linewidth-2.4ex\relax}%
    \raggedright
    \setlength{\parskip}{0.21\baselineskip}\setlength{\parindent}{0pt}
}
\setbeamertemplate{block end}{%
  \vspace{0.9ex}%
  \end{minipage}\hspace*{1.2ex}\par
  \end{beamercolorbox}\vspace{0.20cm}
}

% centre logos on a shared optical line, nudged slightly above the surrounding
% name baseline so the three marks sit comfortably within the header band
\newcommand{\vlogo}[1]{\raisebox{\dimexpr-0.5\height+0.8ex\relax}{#1}}
% centre a photo chip on the optical middle of the adjacent name, half an
% x-height above the baseline
\newcommand{\vphoto}[1]{\raisebox{\dimexpr-0.5\height+0.8ex\relax}{#1}}

% a headshot, or a small labeled placeholder if the file is missing. The
% placeholder parbox is bottom-anchored so it hangs from the baseline like a
% real image would, and \vlogo centers both the same way.
\newlength{\photoside}\setlength{\photoside}{4.5cm}
\newcommand{\photoslot}[2]{% file, label
  \IfFileExists{figures/#1}%
    {\includegraphics[height=\photoside]{#1}}%
    {\fcolorbox{softgray}{white}{\parbox[b][\photoside][c]{\photoside}{\centering
       \textcolor{softgray}{\footnotesize #2}}}}}

% a logo, or a labeled placeholder box if the file is missing
\newcommand{\logoslot}[4]{% file, label, height, placeholder width
  \IfFileExists{figures/#1}%
    {\includegraphics[height=#3]{#1}}%
    {\fcolorbox{softgray}{white}{\parbox[c][#3][c]{#4}{\centering
       \textcolor{softgray}{\small #2\\ {\tiny drop \texttt{\detokenize{figures/#1}}}}}}}}

% Prefer the pinned v8 assets, then fall back to earlier asset sets,
% and show a labeled placeholder rather than a failed build.
\newcommand{\figslot}[3]{% file, width fraction, label
  \IfFileExists{artifacts/figures_v8/#1}%
    {\includegraphics[width=#2\linewidth]{artifacts/figures_v8/#1}}%
    {\IfFileExists{artifacts/figures_v7/#1}%
    {\includegraphics[width=#2\linewidth]{artifacts/figures_v7/#1}}%
    {\IfFileExists{artifacts/figures_v6/#1}%
    {\includegraphics[width=#2\linewidth]{artifacts/figures_v6/#1}}%
    {\IfFileExists{artifacts/figures/#1}%
    {\includegraphics[width=#2\linewidth]{#1}}%
    {\IfFileExists{figures/#1}%
      {\includegraphics[width=#2\linewidth]{#1}}%
      {\fcolorbox{softgray}{white}{\parbox[c][6cm][c]{#2\linewidth}{\centering
         \textcolor{softgray}{\small #3\\ {\tiny missing \texttt{\detokenize{#1}}}}}}}}}}}}

\newcommand{\typicalsectionlead}{\vspace{-0.10cm}}
\newcommand{\typicalsectionfigure}{%
  \figslot{typical_set.pdf}{0.95}{density $\times$ volume and the typical shell}}
\newcommand{\typicalsectiontail}{%
  \vspace{0.08cm}
  Because volume grows as $r^{D-1}$, probability concentrates in the
  \emph{typical set}: a thin shell near $\|\theta\|^2 \approx D$.
  A correct sampler therefore spends almost all its time on this shell.}

\newcommand{\mcmcsectiontitle}{Minibatch gradients are cheap but noisy}
\newcommand{\mcmcsectionbody}{%
  Markov chain Monte Carlo (MCMC) generates a sequence of correlated samples
  targeting $\pi$. Gradient-based samplers use $\nabla\log\pi(\theta)$ to guide
  their updates. Computing this gradient exactly uses all $N$ data points. A
  minibatch of size $B \ll N$ is unbiased and roughly $N/B$ times
  cheaper, but never exact:
  \vspace{-0.18cm}
  \[
    \widehat g_B(\theta)
    = \nabla\log p(\theta)
      + \frac{N}{B}\sum_{i\in\mathcal B}\nabla\log p(\mathcal D_i\mid\theta),
    \qquad \operatorname{Cov}[\widehat g_B] \propto 1/B.
  \]

  \vspace{-0.22cm}
  {\small\textit{Example:} UCI Superconductivity Bayesian linear regression
  ($D=81$, $N=17{,}010$).}

  \vspace{-0.32cm}
  \begin{center}
    \figslot{minibatch_gradient.pdf}{0.96}{minibatch gradients against the exact one}
  \end{center}
  \vspace{-0.35cm}}

\newcommand{\noiseblocktitle}{Noise heats HMC but only turns a ray}
% v7: the mechanism figure gained a third position-space tile, so it takes the
% full figure minipage and a slightly wider share of the row to keep each tile
% readable at distance.
\newcommand{\noisegeometryscale}{1.0}
\newcommand{\noisegeometrypre}{\vspace{-0.45cm}}
\newcommand{\noisefigurewidth}{0.84}
\newcommand{\noisecodewidth}{0.14}
\newcommand{\noiseqrheight}{4.02cm}
\newcommand{\noiseblockopening}{%
  Hamiltonian Monte Carlo (HMC)\refmark{1} and ray tracing (RT)\refmark{2}
  are gradient-based samplers that both use $\nabla\log\pi$. HMC evolves an
  unconstrained velocity, whereas RT evolves a unit direction.

  \textbf{1.\ One noisy step}

  Without a Metropolis correction, mean-zero gradient noise raises HMC's
  expected kinetic energy, increasing its speed.

  For RT, write the isotropic gradient error as $\sigma\zeta$. With
  $n(\theta)=\pi(\theta)^{1/(D-1)}$, the unit direction obeys
  \vspace{-0.03cm}
  \[
      \frac{\dd u}{\dd s}
      = P_u\,\bigg(\nabla \log n(\theta) + \frac{\sigma}{D-1}\,\zeta \bigg),
      \qquad
      P_u = I - uu^\top .
  \]
  \vspace{-0.45cm}

  $P_u$ projects onto the tangent plane, so noise can rotate $u$ but cannot
  change its norm.}

\begin{document}
\begin{frame}[t]

% ---------------- header band ----------------
% top margin matched to the 2.13 cm side margin
\vspace{1.70cm}
\headerband{%
  \rlap{\vlogo{\logoslot{dept_logo_trim.png}{U of T Statistical Sciences signature}{4.3cm}{12cm}}}%
  \parbox[c]{\bandwidth}{\centering
    {\textcolor{uoftblue}{\textbf{\huge Toward Scalable Bayesian Uncertainty for Machine Learning\\[0.3ex] with the Ray Tracing Sampler}}}\\[0.70cm]
    {\LARGE
      \vphoto{\photoslot{photo_yasir_alsugair.png}{Yasir}}\hspace{0.55cm}\AuthorOne{}
      \hspace{2.2cm}
      \vphoto{\photoslot{Chuxuan.png}{Chuxuan}}\hspace{0.55cm}\AuthorTwo{}
      \hspace{3.0cm}
      {\large\textcolor{softgray}{Supervisors:}\ Prof.\ Joshua Speagle and Prof.\ Ricardo Baptista}}}%
  \llap{%
    \vlogo{\logoslot{dsi_logo_standalone.png}{Data Sciences Institute logo}{3.0cm}{12cm}}%
    \hspace{3.0cm}%
    \vlogo{\logoslot{kaust_academy_logo_trim.png}{KAUST Academy logo}{6.0cm}{12cm}}}%
}
\vspace{0.22cm}
\bandrule
\vspace{0.18cm}

% ---------------- three columns ----------------
\begin{columns}[t]

% ============ column 1: why sample, where it lives, how you draw ============
\begin{column}{0.29\paperwidth}

\begin{block}{A trained network is a single guess}
\vspace{0.34cm}

\begin{center}
\figslot{data_bands.pdf}{0.62}{posterior draws vs one Adam fit}
\end{center}
\vspace{-0.3cm}
Ordinary training gives one set of weights, one function. Bayesian sampling
gives a distribution over weights \(\theta\), so uncertainty over functions.
\vspace{-0.05cm}
\[
  \pi(\theta) \;=\; p(\theta \mid \mathcal{D}) \;\propto\;
  \underbrace{p(\mathcal{D} \mid \theta)}_{\text{fit}}\;
  \underbrace{p(\theta)}_{\text{prior}}
\]

\end{block}

\begin{block}{Probability lives on a thin shell}
  \vspace{0.3cm}
\typicalsectionlead
\begin{center}
\typicalsectionfigure
\end{center}
\typicalsectiontail
\end{block}


\begin{block}{\mcmcsectiontitle}
\vspace{0.3cm}
\mcmcsectionbody
\end{block}

\end{column}

% ===================== middle column =====================
\begin{column}{0.31\paperwidth}

\begin{block}{\noiseblocktitle}

\noiseblockopening

% The QR rides in the schematic's spare margin (zero vertical cost): the demo's
% velocity inset is this figure animated, so it is captioned as its live twin.
\noisegeometrypre
\noindent
\begin{minipage}[c]{\noisefigurewidth\linewidth}
  \centering
  \figslot{noise_geometry.pdf}{\noisegeometryscale}{HMC heating versus RT rotation}
\end{minipage}\hfill
\begin{minipage}[c]{\noisecodewidth\linewidth}
  \centering
  \qrcode[height=\noiseqrheight,level=H]{https://chuxuan-a.github.io/rt-demo/}\\[0.5ex]
  \vspace{0.3cm}
  {\small \textbf{Live demo}}
\end{minipage}

\vspace{0.10cm}
\textbf{2.\ Along a trajectory}

With step size \(h\) and trajectory length \(S\), the \(S/h\) noisy kicks
accumulate as
\vspace{-0.12cm}
\[
  \frac{S}{h} \cdot
  O(\sigma^2h^2)
  \;=\; O(\sigma^2h) .
\]
\vspace{-0.45cm}

Let \(\sigma_c(h)\) be the largest noise scale at step \(h\) that keeps long-run
bias below a fixed tolerance, with \(\sigma_c(h)\propto h^{-p}\). For unadjusted
minibatch HMC, \(p=\tfrac12\).

\vspace{0.10cm}
\textbf{3.\ Long-run bias}

For RT, isotropic rotations leave the uniform direction distribution unchanged,
so the \(O(\sigma^2h)\) effect creates no leading bias:
\vspace{-0.05cm}
\[
  \eqbox{\;
  \begin{aligned}
    b_{\mathrm{RT}}(h,\sigma) &= O(\sigma^2h^2)+o(\sigma^2)\\[3pt]
    &\Longrightarrow\;\; \sigma_c(h)=\Omega(h^{-1}),\quad p\ge1
  \end{aligned}
  \;}
\]
\vspace{-0.42cm}
\end{block}

\begin{block}{Measured noise tolerance and sampling cost}
Across \(D=16\)--\(1024\), \(\sigma_c\propto h^{-1}\) for RT and
\(\sigma_c\propto h^{-1/2}\) for HMC; real minibatch noise weakens RT's scaling.

\begin{center}
\makebox[\linewidth][c]{%
  \figslot{tolerance_erosion_frontier.pdf}{1.04}{tolerance, erosion and sampling cost}}
\end{center}
\end{block}

\end{column}

% ============ column 3: the measurement, then the real test ============
\begin{column}{0.31\paperwidth}

\begin{block}{Empirical study on real Gaia data}
We predict each star's alpha abundance plus a per-star scatter from its
Gaia XP spectrum (110 coefficients), for 126{,}156 giants with APOGEE
labels\refmark{4,5}, using a \(D = 11{,}394\) network sampled with
minibatch ray tracing.
% The two panel rows below carry the only real-survey result on the board.
\begin{center}
{\footnotesize Key: \textbf{\textcolor{gaiamap}{point (MAP)}},
\textbf{\textcolor{gaiaensemble}{deep ensemble}},
\textbf{\textcolor{gaiasghmc}{SGHMC}}\refmark{3},
\textbf{\textcolor{gaiart}{ray tracing}}}\\[0.04cm]
\figslot{gaia_calibration_by_group.pdf}{0.78}{per-group calibration, z std per label-error bin}\\[-0.78cm]
{\footnotesize z std = actual / claimed error: closer to 1 is better.}\\[0.52cm]
{\footnotesize averaging the draws beats the single best fit:}\\[0.04cm]
\figslot{exp7_cloud_poster.pdf}{0.76}{one star through 50 posterior draws: tight cloud, trusted prediction}\\[0.06cm]
\figslot{exp7_marginal_poster.pdf}{0.70}{the same two stars in \(y\): MAP spike vs marginalized mass}
\end{center}
\end{block}

\begin{block}{Conclusion}
% Provisional: conclusion wording remains TBD for the final poster.
\begin{itemize}\setlength{\itemsep}{0.02cm}
  \item We took the ray tracing sampler from its paper\refmark{2} to a
    real survey posterior. The draws are the product: averaging them beat
    the single best fit, and their spread warns us when the model is
    guessing.
  \item The gains held on real data: held-out stars were 27\% more
    probable than under the single fit, and an effective sample took
    3.3\(\times\) less data than full-batch sampling.
  \item The theory held as well: gradient noise can only turn a
    constant-speed ray, and the predicted noise tolerance (\(p=1\),
    against \(p=\tfrac12\) for HMC) matched at every dimension we tested.
\end{itemize}
\end{block}

{\footnotesize\setlength{\parskip}{0.20\baselineskip}
\noindent\textcolor{uoftblue}{\rule{\linewidth}{0.7mm}}\\[0.35ex]
\textbf{References}\par
\refitem{1}{Betancourt, M.\ 2017, \emph{A Conceptual Introduction to Hamiltonian Monte Carlo}, arXiv:1701.02434}
\refitem{2}{Behroozi, P.\ 2026, \emph{The Ray Tracing Sampler: Bayesian Sampling of Neural Networks for Everyone}, Open Journal of Astrophysics 9, doi:10.33232/001c.165932}
\refitem{3}{Chen, T., Fox, E.\,B., Guestrin, C.\ 2014, \emph{Stochastic Gradient Hamiltonian Monte Carlo}, ICML, PMLR 32(2), 1683}
\refitem{4}{Gaia Collaboration, Vallenari, A., et al.\ 2023, \emph{Gaia Data Release 3: Summary of the content and survey properties}, A\&A 674, A1}
\refitem{5}{Laroche, A., Speagle, J.\,S.\ 2025, \emph{Closing the Stellar Labels Gap: Stellar Label Independent Evidence for [\(\alpha\)/M] Information in Gaia BP/RP Spectra}, ApJ 979, 5}
\vspace{0.30cm}

\noindent\textbf{Acknowledgments}\ \ Supported by the Data Sciences Institute,
University of Toronto, and the KAUST Academy.
}

\end{column}

\end{columns}

\end{frame}
\end{document}
